\section{Phase mixing as decorrelation of slices}\label{sec:decorrelation}

Without collisions, sources or dissipation of $G$, the energy of every slice is conserved (section~\ref{sec:invariants}). A Hermite moment $g_m$ is a fixed linear combination of slices, so free energy can move between Hermite moments only through changes in the correlations between different slices. Throughout, $C(u,u')=\langle a_u\,a_{u'}\rangle$ is the two-slice correlation, whose diagonal is the slice energy $\mathcal E(u)$ of \eqref{eq:Qphi}, and $W(m)=p_m\langle g_m^2\rangle$ is the Hermite spectrum of the free energy, with $\sum_mW(m)=\int C(u,u)\,\dd u=2W^\pm/n_iT_i$.

\subsection{The two-slice law}\label{dec:sec-twoslice}

We write the slice equation (\ref{eq:slice}) as $\partial_ta_u=-\{\Phi,a_u\}-su\,\gradpar a_u$, using $\partial_z+\{\Psi,\cdot\}=\gradpar$, and take the slice amplitudes to be smooth in $\boldsymbol r$, $t$ and $u$. The bracket and $\gradpar$ obey the product rule, so for two slices
\begin{equation*}
  \partial_t(a_ua_{u'})=-\{\Phi,a_ua_{u'}\}-su'\,\gradpar(a_ua_{u'})-s\,(u-u')\,a_{u'}\gradpar a_u .
\end{equation*}
Since $\nabla\cdot\uE=\nabla\cdot\bhat=0$, $\{\Phi,f\}=\nabla\cdot(f\,\uE)$ and $\gradpar f=\nabla\cdot(f\,\bhat)$, so the first two terms on the right are divergences, and their averages over the periodic box vanish. The two-slice correlation therefore obeys
\begin{equation}\label{eq:twoslice}
  \frac{\dd C(u,u')}{\dd t}=-s\,(u-u')\,\big\langle a_{u'}\,\gradpar a_u\big\rangle .
\end{equation}
In particular, $C(u,u)$ is constant for every $u$, and the $\boldsymbol E\times\boldsymbol B$ flow does not enter.

The $\boldsymbol E\times\boldsymbol B$ flow carries the two slices of a pair together and cannot change their correlation, which changes only through their relative velocity $s(u-u')\,\bhat$, that is, through relative streaming along the perturbed field. Phase mixing and phase unmixing are therefore the decorrelation and the recorrelation of slices along $\gradpar$, a definition that refers to no wavenumber. When the induction equation is ideal and the Alfv\'enic fields stay smooth, every volume-integrated single-point statistic of $G$, the correlations $C(u,u')$ among them, takes the value it has in linear phase mixing along the field lines frozen at $t=0$ \citep{KanekarPrep}.

The law needs no assumption on the Alfv\'enic fields beyond smoothness; terms that act on $G$ add those of \eqref{eq:sliceforcing}. It holds for the truncated hierarchy, with $C_{jl}=\langle a_ja_l\rangle$ at the nodes $u_j$, under the conditions of section~\ref{sec:truncated}, and we checked it to round-off against GANDALF's right-hand side \citep{Kanekar2025}.

With a straight guide field, in electrostatic drift kinetics \citep{Schekochihin2016,Parker2016} and in the linear model of \citet{Kanekar2015}, the Hermite moments of each Fourier mode are split by the sign of the parallel wavenumber into parts that propagate to larger and to smaller $m$. In KRMHD the field lines move; \citet{Meyrand2019} noted that the sign would then have to be taken with respect to the perturbed field and did not use such a split. The two-slice law describes phase mixing along the perturbed field without reference to a wavenumber. Where the field lines are tilted, a split by the sign of $k_z$ exchanges phase mixing and unmixing on modes whose $k_z$ and field-aligned wavenumber have opposite signs, whereas with an ideal induction equation no part of the distribution changes the sign of its field-aligned wavenumber \citep{KanekarPrep}.

\subsection{The Hermite flux as a Christoffel--Darboux transform}\label{dec:sec-flux}

With the orthonormal functions $\phi_n=\sqrt{p_n\varrho}\;r_n$ of section~\ref{sl:forcingsec}, (\ref{sl:parseval}) gives
\begin{equation}\label{dec:gn}
\begin{gathered}
  g_n=\frac{1}{\sqrt{p_n}}\int\phi_n(u)\,a_u\,\dd u,\\
  E_{\le m}\equiv\sum_{n\le m}W(n)=\iint K_m(u,u')\,C(u,u')\,\dd u\,\dd u',\qquad
  K_m(u,u')=\sum_{n\le m}\phi_n(u)\,\phi_n(u') .
\end{gathered}
\end{equation}
Let $b_m=p_mA_{m,m+1}=p_{m+1}A_{m+1,m}$. We multiply the equation for $g_n$ in (\ref{mod:hierarchy}) by $2p_ng_n$, average over the box and sum over $n\le m$. The advection drops out because the flow is incompressible, and the streaming couplings cancel in pairs, because $p_nA_{n,n+1}=p_{n+1}A_{n+1,n}$ and $\langle f\gradpar h\rangle=-\langle h\gradpar f\rangle$, except the one between $m$ and $m+1$. The slice form of $E_{\le m}$ with the two-slice law gives the same rate:
\begin{equation}\label{eq:fluxCD}
\begin{aligned}
  \frac{\dd E_{\le m}}{\dd t}=-\Pi_{m+1/2},\qquad
  \Pi_{m+1/2}&=s\,b_m\big\langle g_m\gradpar g_{m+1}-g_{m+1}\gradpar g_m\big\rangle\\
  &=-\iint K_m(u,u')\,\frac{\partial C(u,u')}{\partial t}\,\dd u\,\dd u' ,
\end{aligned}
\end{equation}
the free energy per unit time and volume that passes from moment $m$ to $m+1$. The first form is the Hermite flux of \citet{Schekochihin2016} and \citet{AdkinsSchekochihin2018} with the full field-aligned gradient, as used by \citet[eq.~9]{Meyrand2019}.

The two forms are related by a Christoffel--Darboux identity. Since $\Tr[LG]=u\,\Tr[G]$ and $Le_n=A_{n-1,n}e_{n-1}+A_{n+1,n}e_{n+1}$, the symmetry of $\Pw A$ gives the recurrence $u\,r_n=A_{n,n+1}\,r_{n+1}+A_{n,n-1}\,r_{n-1}$, $r_{-1}=0$. Multiply it at $u$ by $p_nr_n(u')$, subtract the same with $u$ and $u'$ exchanged, and sum over $n\le m$; since $p_nA_{n,n+1}=b_n$ and $p_nA_{n,n-1}=b_{n-1}$, the sum telescopes to
\begin{equation}\label{dec:CD}
  (u-u')\,K_m(u,u')=b_m\sqrt{\varrho(u)\varrho(u')}\,\big[r_{m+1}(u)\,r_m(u')-r_m(u)\,r_{m+1}(u')\big] .
\end{equation}
With \eqref{eq:twoslice}, the second form of $\Pi_{m+1/2}$ is $s\iint(u-u')K_m\langle a_{u'}\gradpar a_u\rangle\,\dd u\,\dd u'$, which \eqref{dec:CD} and \eqref{dec:gn} turn into the first. The pairs $u=u'$ carry no flux: the free energy that passes through every $m$ is carried by the correlations between different slices. The truncated hierarchy has the same structure: for $m<M$ and both closures, rows $0$ to $m$ of $A\boldsymbol r_j=u_j\boldsymbol r_j$ are the recurrence of the $r_n$, so the same telescoping gives $(u_j-u_l)K_m(j,l)=b_m(R_{m+1,j}R_{ml}-R_{mj}R_{m+1,l})$, with $K_m(j,l)=\sum_{n\le m}p_nR_{nj}R_{nl}$ and $R_{mj}$ of (\ref{tr:aj}), normalised with the weight of the closure used. We checked this to round-off with GANDALF's ladder matrix and right-hand side.

\subsection{The Hermite spectrum and the two-slice correlation}\label{dec:sec-dictionary}

By \eqref{dec:gn}, the Hermite spectrum is a quadratic form in the two-slice correlation:
\begin{equation}\label{eq:dictionary}
  W(m)=\iint C(u,u')\,\phi_m(u)\,\phi_m(u')\,\dd u\,\dd u' .
\end{equation}
We call the slices $\delta$-correlated if $C(u,u')=e(u)\,\delta(u-u')$ with a density $e\ge0$. They are then completely decorrelated in velocity, and $W(m)=\int e\,\phi_m^2\,\dd u$. For $e=F_M$ at $\Lam=\infty$,
\begin{equation}\label{eq:prefactor}
  \int F_M\,\psi_m^2\,\dd u=\frac{\Gamma(m+\frac12)}{\sqrt2\,\pi\,m!}=\frac{1}{\pi\sqrt{2m}}\Big(1-\frac{1}{8m}+O(m^{-2})\Big),
\end{equation}
where $\Gamma(\cdot)$ is Euler's gamma function. This spectrum decays only as $m^{-1/2}$, so $\sum_mW(m)$ diverges: $\delta$-correlated slices idealise states with velocity structure on arbitrarily fine scales, such as the collisionless forced state below.

\citet{Kanekar2015} studied a one-dimensional linear model of the form of one branch, with $\Lam=-1/\alpha$ for their parameter $\alpha$, whose density moment is driven by white noise of strength $\epsilon_k$. Their eq.~(4.10), derived without collisions, gives the exact steady-state Hermite spectrum of a Fourier mode for $m\ge1$, which in our variables (their $|k|$ is our $s|k_z|$) reads
\begin{equation*}
  W(m)=\frac{(1-1/\Lam)^2\,\epsilon_k}{2\pi\,s|k_z|}\,\frac{1}{2^mm!}\int\frac{|Z^{(m)}(\zeta)|^2}{|\diel(\zeta+\ii0)|^2}\,\dd\zeta ,
\end{equation*}
where $Z^{(m)}$ is the $m$th derivative of the plasma dispersion function. At $\Lam=\infty$, $\diel=1$, and on the real axis $\mathrm{Im}\,Z^{(m)}=(-1)^m\pi H_mF_M$ and $\mathrm{Re}\,Z^{(m)}=(-1)^m\pi\,\mathsf H[H_mF_M]$, with $\mathsf H$ the Hilbert transform, which preserves $\int f^2\,\dd u$. Hence $(2^mm!)^{-1}\int|Z^{(m)}|^2\,\dd\zeta=\sqrt2\,\pi\,\Gamma(m+\frac12)/m!$, and
\begin{equation}\label{dec:kanekar}
  W(m)=\frac{\pi\epsilon_k}{s|k_z|}\,\frac{\Gamma(m+\frac12)}{\sqrt2\,\pi\,m!} ,
\end{equation}
which is \eqref{eq:prefactor} with $e=(\pi\epsilon_k/s|k_z|)F_M$: the forced linear spectrum is the spectrum of $\delta$-correlated slices. The closed form is not written in their paper, whose eq.~(4.14) is its large-$m$ limit, and the slice argument covers $m=0$ as well.

The forced slices are $\delta$-correlated because they dephase: density forcing loads every slice with the same spatial pattern \eqref{eq:sliceforcing}, two slices loaded together dephase at the rate $s|k_z(u-u')|$, and the steady state superposes all past injections, so the real part of the mean of $a_u(\boldsymbol k)\,a_{u'}(\boldsymbol k)^*$ vanishes for $u\ne u'$. The antisymmetric imaginary part, which does not enter \eqref{eq:dictionary}, carries the injection rate $\epsilon_W=(1-1/\Lam)\,\epsilon_k$ through every $m$. At finite $\Lam$ the slices are again $\delta$-correlated in the mean, with $e=\pi(1-1/\Lam)^2\epsilon_k\,\varrho/(s|k_z|)$, which reproduces their eq.~(4.10) exactly for every $m\ge1$. Proportionality to $(1-1/\Lam)$ times \eqref{dec:kanekar} holds only at large $m$, where both tend to $\epsilon_W/(s|k_z|\sqrt{2m})$, independent of $\Lam$ at fixed $\epsilon_W$ by their eq.~(4.15), which in our variables is $\int\varrho\,\dd u=1/p_0$.

That $\sqrt{2m}$ acts as a wavenumber in velocity space is the duality noted by \citet{Schekochihin2016} and \citet{AdkinsSchekochihin2018}; \citet[note~81]{Eyink2018} related the $m^{-1/2}$ spectrum to a velocity-space regularity exponent of $-1/2$, and \citet{Nastac2024} describe linear phase mixing as producing a spectrum flat in the velocity-space wavenumber. The slice picture adds the representation \eqref{eq:dictionary}, the closed form \eqref{dec:kanekar}, and the reading of the $m^{-1/2}$ spectrum as complete decorrelation of the slices in velocity.

\subsection{No Fj{\o}rtoft-type constraint}\label{dec:sec-fjortoft}

\citet[section~4.1]{Chandran2026} discuss cascade directions in $m$ with the generalisation by \citet{AlexakisBiferale2018} of the argument of \citet{Fjortoft1953}: if $E_{\mathcal I_2}(m)\ge c\,m^n|E_{\mathcal I_1}(m)|$ for two invariants and positive $c$ and $n$, then $\mathcal I_1$ cannot cascade to large $m$ at a constant rate. In their appendix~B they note that this condition fails for $W$ and $\Gamma$, which is not sign-definite, and expect both to cascade to large $m$. We complete their argument for $\Gamma$ and extend it to every invariant that depends on $G$ alone. Let $\Gamma_m=2b_m\langle g_mg_{m+1}\rangle$, so that $\sum_m\Gamma_m=\langle\fe{G}{LG}\rangle$; up to normalisation, $\Gamma_m$ is the contribution of the pair $(m,m+1)$ to CMM's (B32).

The simplest solutions of collisionless KRMHD already rule out such a constraint. With $\Phi=\Psi=0$, which solve \eqref{eq:rmhd}, the slice equation \eqref{eq:slice} reduces to $\partial_ta_u+su\,\partial_za_u=0$, so every $\Qphi$ and every Casimir \eqref{eq:casimir} is conserved, as is every other invariant of collisionless KRMHD that depends on $G$ alone. Consider one branch, with $\Lam>1$, $\Lam<0$ or $\Lam=\infty$, in a periodic box, and an initial state that is smooth in space, has a finite Hermite expansion in $u$ and positive free energy $\langle\fe{G}{G}\rangle$, and has no Fourier component at $k_z=0$. By \eqref{dec:gn}, the Hermite moments of each Fourier mode are then
\begin{equation}\label{dec:freestream}
  g_m(\boldsymbol k,t)=\int\sqrt\varrho\,r_m\,a_u(\boldsymbol k,0)\,e^{-\ii sk_zut}\,\dd u ,
\end{equation}
integrals over $u$ of a fixed function times a phase that oscillates ever faster in $u$. They tend to zero as $t\to\infty$, and so do $W(m,t)$ and $\Gamma_m(t)$ for every $m$.\footnote{The integrand of \eqref{dec:freestream} is integrable in $u$ by the Cauchy--Schwarz inequality, so $g_m(\boldsymbol k,t)\to0$ by the Riemann--Lebesgue lemma. Since $|g_m(\boldsymbol k,t)|^2\le p_m^{-1}\int|a_u(\boldsymbol k,0)|^2\,\dd u$, which is summable over $\boldsymbol k$, the sums over $\boldsymbol k$ that give $W(m,t)$ and $\Gamma_m(t)$ tend to zero as well.} The free energy, which is conserved, is carried to ever larger $m$, while every invariant that depends on $G$ alone stays constant.

Suppose now that a quantity $\mathcal I$ that depends on $G$ alone and is conserved by these solutions, such as any invariant of collisionless KRMHD that depends on $G$ alone, could be split over Hermite indices, $\mathcal I=\sum_mE_{\mathcal I}(m)$, in such a way that $E_{\mathcal I}(m)\ge c\,m^nW(m)$ for all $m$, whatever the state, with constants $c>0$ and $n>0$, as the enstrophy and energy spectra of two-dimensional flow obey $Z(k)=k^2E(k)$. Then $E_{\mathcal I}(m)\ge0$, and $\mathcal I\ge c\,m_*^n\sum_{m>m_*}W(m,t)$ for every $m_*$. As $t\to\infty$ the sum tends to the whole free energy $\langle\fe{G}{G}\rangle$, so for $m_*$ large enough the right-hand side exceeds the constant $\mathcal I$, which is a contradiction. The same argument, started from a state with $\langle\fe{G}{LG}\rangle\neq0$, excludes $E_{\mathcal I}(m)\ge c\,m^n|\Gamma_m|$: since $uG$ keeps a finite free energy, $\sum_m\Gamma_m(t)$ converges to the conserved $\langle\fe{G}{LG}\rangle$, so that $\sum_{m>m_*}|\Gamma_m(t)|\ge|\langle\fe{G}{LG}\rangle-\sum_{m\le m_*}\Gamma_m(t)|\to|\langle\fe{G}{LG}\rangle|$.

Moreover, every $\Gamma_m$, like $\Gamma$ and $\Hph$ (section~\ref{sl:special}), is odd under the mirror map \eqref{eq:mirror}. Since $H_m(-u)=(-1)^mH_m(u)$, the map takes $g_m$ to $(-1)^mg_m(\tilde{\boldsymbol r})$, so it reverses the sign of every $\langle g_mg_{m+1}\rangle$ and leaves $W(m)$ unchanged. A mirror-symmetric state therefore has $\Gamma_m=0$ for all $m$ and $\Hph=0$, whatever its $W(m)$.

Linear phase mixing in a uniform field thus conserves every invariant of section~\ref{sec:invariants} while it moves all of the free energy, and all of $\Gamma$, to arbitrarily large $m$ (in a truncated hierarchy, until the recurrence time); in two-dimensional hydrodynamics this is impossible, because the energy above a wavenumber $k_*$ is bounded by the enstrophy divided by $k_*^2$. The reason is that every invariant of section~\ref{sec:invariants} depends only on statistics of slices at one $u$, whereas the Hermite transfer \eqref{eq:fluxCD} is $-\iint K_m\,\partial_tC$, a change in the correlations between slices at different $u$. The expectation of \citet{Chandran2026} that $\Gamma$, like $W$, is transferred to large $m$ is therefore borne out by these exact solutions, for a reason that does not rely on the sign-indefiniteness of $\Gamma$. The argument concerns transient solutions of collisionless KRMHD; CMM's criterion refers to inertial-range fluxes in a forced and dissipative steady state, which this argument does not address.
